%% General definitions
%\documentclass{article} %% Determines the general format.
\documentclass[10pt,a4paper]{article}
\usepackage{a4wide} %% paper size: A4.
\usepackage[utf8]{inputenc} %% This file is written in UTF-8.
%% Some editors on Windows cannot save files in UTF-8.
%% If there is a problem with special characters not showing up
%% correctly, try switching "utf8" to "latin1" (ISO 8859-1).
\usepackage[T1]{fontenc} %% Format of hte resulting PDF file.
\usepackage{fancyhdr} %% Package to create a header on each page.
\usepackage{lastpage} %% Used for "Page X of Y" in the header.
%% For this to work, you have to call pdflatex twice.
\usepackage{enumitem} %% Used to change the style of enumerations (see below).

\usepackage{amssymb} %% Definitions for math symbols.
\usepackage{amsmath} %% Definitions for math symbols.
\usepackage{amsthm}
\usepackage{amsfonts}
\usepackage{mathtools}

\usepackage[english]{babel}

\usepackage{graphicx}
\usepackage{hyperref}
\setlength{\parindent}{0cm}

%\usepackage[left=2.00cm, right=2.00cm, top=2.00cm, bottom=2.00cm ]{geometry}

\newtheorem*{zz*}{t.s}
\newtheorem*{theorem*}{Theorem}
\newtheorem*{claim*}{claim}
\newtheorem{claim}{claim}
\newtheorem*{prop*}{Propsition}
\newtheorem*{lemma*}{Lemma}
\newtheorem{lemma}{Lemma}

%\title{Lineare Algebra Serie 7}
\date{\today}


\usepackage{tikz}  %% Pagacke to create graphics (graphs, automata, etc.)
\usepackage{tikz-cd} %%commutativ system
\usetikzlibrary{automata} %% Tikz library to draw automata
\usetikzlibrary{arrows}   %% Tikz library for nicer arrow heads


%% Left side of header
\lhead{\course\\\semester\\Serie \homeworkNumber}
%% Right side of header
\rhead{\authorname\\Page \thepage\ of \pageref{LastPage}}
%% Height of header
\usepackage[headheight=36pt]{geometry}
%% Page style that uses the header
\pagestyle{fancy}

\newcommand{\authorname}{Ken Walzer}
\newcommand{\semester}{HS 2023}
\newcommand{\course}{Algebraic number theory}
\newcommand{\homeworkNumber}{1}
\newcommand{\Zn}{\mathcal{Z}}
\newcommand{\Vn}{V_{\mathbb{P}^n}}
\newcommand{\RR}{\mathbb{R}}
\newcommand{\real}{\mathrm{Re}}
\newcommand{\imag}{\mathrm{Im}}
\newcommand{\CC}{\mathbb{C}}
\newcommand{\CCC}{\mathcal{C}}
\newcommand{\QQ}{\mathbb{Q}}
\newcommand{\ZZ}{\mathbb{Z}}
\newcommand{\NN}{\mathbb{N}}
\newcommand{\FF}{\mathbb{F}}
\newcommand{\Mat}{\mathrm{Mat}}
\newcommand{\Aut}{\mathrm{Aut}}
\newcommand{\PP}{\mathbb{P}}
\newcommand{\A}{\mathbb{A}}
\newcommand{\OO}{\mathcal{O}}
\newcommand{\kk}{\mathbf{k}}
\newcommand{\aId}{\mathfrak{a}}
\newcommand{\ii}{\sqrt{-1}}
\newcommand{\GL}{\mathrm{GL}}
\newcommand{\SL}{\mathrm{SL}}
\newcommand{\orb}{\mathrm{orb}}
\newcommand{\stab}{\mathrm{stab}}
\newcommand{\ord}{\mathrm{ord}}
\newcommand{\id}{\mathrm{id}}
\linespread{1.2}

\begin{document}
	
	\section*{Aufgabe 1}
	\begin{zz*}
		It holds that:
		\begin{enumerate}[label = (\roman*)]
			\item $\ZZ[i]\subseteq \CC$ is a subring
			\item $\ZZ[i]$ forms a free abelian group with rank $2.$ and $\ZZ[i] = \ZZ+i\ZZ.$
			\item Let $u\in \ZZ[i]^*$ be a unit. It holds that $|u| = 1.$
			\item $\ZZ[i]^* = \{\pm1,\pm i\}.$
			\item $5 = (2+i)(2-i)$
			\item Let $a,b\in \ZZ[i],$ such that $3 = ab.$ Then $a$ or $b$ is a unit. In particular $3$ is irreducible in $\ZZ[i].$
			\item $\ZZ[i]\cap \QQ = \ZZ.$
		\end{enumerate}
	\end{zz*}
	Before we proof the above statement we show a small Lemma.
	\begin{lemma*}
		It holds that $\ZZ[i] = \ZZ+i\ZZ.$
	\end{lemma*}
	\begin{proof}
		Let $a\in\ZZ[i].$ Then there exists a $f = \sum_{i=0}^nf_iX^i\in\ZZ[X],$ such that $a = f(i).$Then  denote $N = \{0,\dots, n\}.$ Furthermore we define 
		\[T_0 = 4\ZZ\cap N,\quad T_1 = 1+4\ZZ\cap N,\quad T_2 = 2+4\ZZ\cap N,\quad T_3 = 3+\ZZ\cap N. \]
		It holds that  $N = T_0 +T_1+T_2+T_3.$ Then we define for each $i\in \{1,2,3,4\},$ the number $F_i = \sum_{k\in T_i} f_k\in \ZZ.$ This gives us, that 
		\[\begin{split}
			a &= f(i) = \sum_{k=0}^nf_ki^k = \sum_{k\in T_0}^nf_ki^k + \sum_{k\in 1+4\ZZ} f_k i^k +\sum_{k\in 2+4\ZZ} f_k i^k + \sum_{k\in 3+4\ZZ} f_k i^k\\
			&= 1\cdot \left(\sum_{k\in T_0}^nf_k\right) + (i)\cdot \left(\sum_{k\in T_1}^nf_k\right) + (-1)\cdot \left(\sum_{k\in T_2}^nf_k\right) + (-i)\cdot \left(\sum_{k\in T_3}^nf_k\right)\\
			&= F_0-F_2 + i\cdot (F_1-F_3)
		\end{split}\]
		It follows that $a\in \ZZ+i\ZZ$ and therefore, that $\ZZ[i]\subseteq \ZZ+i\ZZ.$ On the otherhand let $a\in \ZZ + i\ZZ.$ Then there exists $a_1,a_2\in \ZZ,$ such that $a = a_1+a_2i.$ But this concides with the evaluation of the polynomial $a_1+a_2X\in \ZZ[X]$ at $i.$ So the other inclusion follows and therefore its equality and we can conclude.
	\end{proof}
	Now we want to prove the exercise.
	\begin{proof}
		\begin{enumerate}[label = (\roman*)]
			\item Let $a,b\in \ZZ[i].$ By above Lemma there exist $a_1,a_2,b_1,b_2\in \ZZ,$ such that $a = a_1+a_2i, b = b_1+b_2i.$ Then it holds that $\ZZ[i]\neq \emptyset,$ since $\ZZ\subseteq \ZZ[i]$ and
			\[\begin{split}
				a-b &= a_1-b_1 + (a_2 - b_2)i \in \ZZ[i]\\
				a\cdot b &= (a_1b_1 - a_2b_2) + (a_1b_2 + a_2b_1)i\in \ZZ[i]
			\end{split}\]
			Moreover it holds that $\ZZ[i]\subseteq \CC = \RR(i)$ concluding that $\ZZ[i]\subseteq \CC$ forms a subring.
			\item By above Lemma it holds that $\ZZ[i]$ is generated by $(1, i)$ as $\ZZ$-module. Let $a_1,a_2\in \ZZ,$ satisfying that
			\[a_1 + a_2i = 0\]
			This implies that $a_1 = -a_2i.$ Hence, $a_1,a_2 = 0$ otherwise $i\in \QQ$ which would imply $\CC = \RR(i) = \RR$ which is absurd. It follows that $(1,i)$ forms a basis of $\ZZ[i].$ Then $\ZZ$ being abelian implies that $\ZZ[X]$ is abelian and therefore its homomorphic image $\ZZ[i]$ is abelian as well. It follows that $\ZZ[i]$ forms a free abelian group with rank $2.$ the second statement is shown in above Lemma.
			\item Let $u\in \ZZ[i]^*.$ Then there exists a $a\in \ZZ[i]^*,$ such that $1 = ua.$ We note that the absolute value on $\CC$ forms a norm by definition 1.5 from script, furthermore is it is mulitplicative since Lemma 1.7 implies that any norm on a finite field extension is a product of field homomorphisms. Together with $x\mapsto x^2$ being an endomorphism, we get that
			\[1 = (|ua|)^2= |u|^2|a|^2 = \underbrace{(\real(u)^2 + \imag(u)^2)}_{\in \ZZ}\underbrace{(\real(a)^2+\imag(a)^2)}_{\in \ZZ}\]
			Since $\pm 1$ the only units in $\ZZ,$ this implies that $|u|^2 = \pm1$ by definition of the absolute value, it follows that $|u|^2=1,$ thus $|u| = 1.$
			\item We consider the embedding $\ZZ[i]\simeq \ZZ^2\subseteq \RR\simeq \RR^2\simeq \CC.$ Which we can do since we have shown that $\ZZ[i]$ is a free $\ZZ-$ module of rank $2.$ (and structure theorem of finitely generated $\ZZ$-modules gives us, that $\ZZ[i]$ forms a lattice of dimension $2$) Then we know by (iii) that $\ZZ[i]^*\subseteq \ZZ^2\cap\partial B_{1}(0) = \{pm1, \pm i\}.$ Furthermore it holds that $(-1)^2, i(-i) = 1,$ thus equality follows. 
			\item It holds that
			\[(2+i)(2-i) = (4+1) + i\cdot 0 = 5\]
			\item By the multiplicativity of the norm, it holds that
			\[3 = |3| = |ab| = |a||b|\in \ZZ\]
			(note that any norm in sense of the script is mutliplicative since a Lemma gives us that the norm can be written as a product of fieldhomomorphisms)
			But we know that $3$ is a prime number and therefore irreducible. It follows that $N(a)$ or $N(b)$ is an unit.
			\item Since $\ZZ\subseteq \ZZ[i], \QQ,$ the iclusion $\ZZ\subseteq \ZZ[i]\cap \QQ$ holds. Let $x\in \ZZ[i]\cap \QQ.$ Then there exists $a,b\in \ZZ,$ such that $x = a+bi.$ Then $bi = x-a\in \QQ.$ If $b\neq 0,$ this would imply $i\in \QQ$ which is absurd. Hence $x = a\in \ZZ$ The other inclusion follows which gives us the equality.
		\end{enumerate}
	\end{proof}
	\pagebreak
	%%%%%%%%%%%%%%%%%%%%%%%%%%%%%%%%%%%%%%%%%%%%%%%%%%%%%%%%%
	\section*{Aufgabe 2}
	\begin{zz*}
		A Ring $R$ is noetherian, if and only if any ideal is finitely generated.
	\end{zz*}
	\begin{proof}
		Assuming $R$ is noetherian. Then let $I\subseteq R$ be an ideal. The case $I = (0)$ is eivdent, so we may assume that $I\neq (0).$ Let $a_0\in I.$ Then $I_0 = aR\subseteq I.$ Then we define inductively for each positive integer $i\in \ZZ,$ an element $a_i\in R$ by
		\[(a_i, I_{i})= \left(\begin{cases}
			\text{some element in } I\setminus I_{i-1} & \text{ if } I_{i-1}\neq I\\
			a_{i-1} & \text{ if } I = I_{i-1}
		\end{cases}, (a_1R+\cdots a_{i}R)\right)\]
		(We note here that there is no need for AOC since we only need finitely many elements as a result)
		
		Then $I_0\subseteq I_{1}\subseteq \cdots$ is a nondecreasing sequence of ideal. By assumpiton, there exists a $n\geq 1$ such that $I_{n} = I_{n+k}$ for all $k\in \NN.$ Resulting in, that $I = (a_0R+\cdots anR)$ which gives us one direction.
		
		Now let's assume that every ideal is finitely generated. Let $(I_n)_n\subseteq \mathcal{P}(R)$ be a nondecreasing sequence on ideals. Then we define
		\[I = \bigcup_{k=0}^\infty I_k.\]
		We claim that $I\subseteq R$ forms an ideal. Indeed , first we get $\emptyset\neq I_0\subseteq I.$ Then let $a,b\in I$ and $r\in R.$ There exists $k,l\in \NN$ such that $a\in I_k,b\in I_l.$ Without loss of generality we may assume that $k\geq l,$ and we get that $a-b,ra\in I_k\subseteq I$ by $I_k$ being an ideal. It follows that $I\subseteq R$ forms an ideal. By assumption, there exist $a_1,\dots, a_n\in I,$ that generate $I.$ Then for each $i\in \{1,\dots, n\},$ there exists a $k_i\in \NN,$ such that $a_i\in I_{k_i}.$
	\end{proof}
	\pagebreak
	%%%%%%%%%%%%%%%%%%%%%%%%%%%%%%%%%%%%%%%%%%%%%%%%%%%%%%%%%%%%%
	\section*{Aufgabe 3}
	\subsection*{(i)}
	Let $I\subseteq \ZZ[i]$ be a non-zero ideal.
	\begin{zz*}
		It holds that $I$ is a free $\ZZ$-module of rank $2.$
	\end{zz*}
	
	\subsection*{(ii)}
	\begin{zz*}
		It holds that $\ZZ[i]$ is noetherian.
	\end{zz*}
	\pagebreak
\end{document}
